\documentclass[reprint, amsmath, amssymb, aps, prl]{revtex4-2}

\usepackage{graphicx} % Include figure files
\usepackage{dcolumn}  % Align table columns on decimal point
\usepackage{bm}       % bold math
\usepackage[hidelinks,hypertexnames=false]{hyperref}

\newcommand{\Qdot}{\dot Q}
\newcommand{\degC}{\ensuremath{\,^\circ\mathrm{C}}}
\newcommand{\Biperp}{\mathrm{Bi}_\perp}
\newcommand{\WmK}{\mathrm{W\,m^{-1}K^{-1}}}
\newcommand{\WmmK}{\mathrm{W\,m^{-2}K^{-1}}}
\urldef{\repourl}\url{https://github.com/tyin-source/__Lab_Module_7__}
\urldef{\pathPone}\path{python/analysis/a4_problem_1_3.py}

\begin{document}

\preprint{Phys 39 Module 7 -- A4}

\title{Finite-Rod Model and Validity Checks:\\ Fin Equation, Finite-Length Error, and Transverse Biot Number}

\author{Tianren Yin}
\affiliation{Physics Department, Brandeis University}

\date{\today}

\begin{abstract}
This assignment builds the stationary finite-rod (fin) model used for the Module 7
course rod and makes two provisional validity predictions before heating. Applied to
Lienhard Problem~1.3, Fourier's law gives temperature drops from $0.1$~K to $2000$~K
across a 1-cm slab at the same heat flux; the three low-conductivity cases are physically
unreasonable. An energy balance on a rod slice gives the fin equation
$d^2\theta/dx^2=m^2\theta$ with $m^2=hP/(kA)$, whose insulated-tip solution is
$\cosh[m(L-x)]/\cosh mL$. For the 6061-T6 rod ($d=12.70$~mm, $L=26.90$~cm,
$k=166$--$172~\WmK$) and the planning value $h=10~\WmmK$, $m=4.34~\mathrm{m^{-1}}$ and
$mL=1.17$, so the semi-infinite profile $e^{-mx}$ is in error by 11--30\% at the three
farthest planned sensors: the rod is not expected to be semi-infinite over the measured
region. In contrast, $\Biperp=hd/(4k)=1.9\times10^{-4}\ll0.1$, so each cross section can
be represented by one temperature.
\\[4pt]
Assignment A4 (individual). Repository: \repourl; Git commit \texttt{\commit}.
\end{abstract}

\maketitle

Parts~3 and 4 are provisional predictions recorded before heating the rod. They use the
measured rod geometry, the planned sensor positions, a literature range for $k$, and the
planning value $h=10~\WmmK$ given in the Part~III handout~\cite{handout}. Equation numbers
in parentheses such as (4.28) are those of Lienhard~\cite{Lienhard}.

%======================================================================
\section*{Part 1. Apply Fourier's law: Lienhard Problem 1.3}

\subsection*{Problem statement}

A flat slab of thickness $L=0.01$~m carries a steady heat flux $q=5000~\mathrm{W/m^2}$
from its hot side to its cold side. The cold side is held at $T_c=-40\degC$. For each
material in Table~\ref{tab:p13}, compute the temperature drop $\Delta T$ across the slab,
find the hot-side temperature $T_h=T_c+\Delta T$, and state whether the case is
physically reasonable and why.

\subsection*{Derivation from Fourier's law}

Start from the one-dimensional Fourier law,
\begin{equation*}
q_x=-k\frac{dT}{dx}.
\end{equation*}
Assumptions:
\begin{enumerate}
\item \textbf{Steady state}: no time dependence, $\partial T/\partial t=0$;
\item \textbf{One-dimensional conduction}: $T=T(x)$ only;
\item \textbf{Constant $k$}: thermal conductivity independent of temperature over the range;
\item \textbf{No internal heat generation}.
\end{enumerate}
With no source and steady state, the energy balance on a slice gives
\begin{equation*}
\frac{dq_x}{dx}=0\quad\Longrightarrow\quad q_x=q=\text{constant}.
\end{equation*}
Then
\begin{equation*}
\frac{dT}{dx}=-\frac{q}{k}=\text{constant},
\end{equation*}
so the temperature profile is \textbf{linear in $x$}. This linearity is a \emph{result} of
the steady, source-free, constant-$k$ assumptions, not an extra assumption. Integrating
from $x=0$ (hot side) to $x=L$ (cold side),
\begin{equation*}
T(L)-T(0)=-\frac{qL}{k}.
\end{equation*}
Defining $\Delta T=T_h-T_c=T(0)-T(L)$,
\begin{equation*}
\boxed{\Delta T=\frac{qL}{k}}
\qquad
\boxed{T_h=T_c+\frac{qL}{k}}.
\end{equation*}

\subsection*{Unit check}

\begin{equation*}
\frac{\mathrm{W/m^2}\cdot\mathrm{m}}{\mathrm{W/(m\cdot K)}}
=\frac{\mathrm{W/m}}{\mathrm{W/(m\cdot K)}}
=\mathrm{K}.
\end{equation*}
A temperature difference in kelvin equals the same difference in $^\circ$C, so $T_h$ can
be reported in $^\circ$C.

\subsection*{Worked example: silver}

For silver, $k=427~\WmK$. With $q=5000~\mathrm{W/m^2}$ and $L=0.01$~m,
\begin{align*}
\Delta T
&=\frac{(5000~\mathrm{W/m^2})(0.01~\mathrm{m})}{427~\mathrm{W/(m\cdot K)}}
=\frac{50}{427}~\mathrm{K}\\
&\approx0.117~\mathrm{K}.
\end{align*}
Hence
\begin{equation*}
T_h=-40\degC+0.117~\mathrm{K}\approx-39.88\degC.
\end{equation*}
The hot side stays far below silver's melting point, and the temperature range is tiny,
so constant-$k$, steady, one-dimensional conduction is an excellent approximation.

\subsection*{Table calculation for all materials}

For every material the same formula applies,
\begin{equation*}
\Delta T=\frac{qL}{k}=\frac{50}{k}~\mathrm{K},
\qquad
T_h=-40\degC+\Delta T,
\end{equation*}
with the results in Table~\ref{tab:p13}. The metals use the 20\degC\ column, as the
book's silver answer does. With $k$ interpolated to $-40\degC$, silver gives 0.1165~K
and aluminum 0.2097~K, a change below 1\%.

\begin{table*}[t]
\caption{\label{tab:p13}Lienhard Problem 1.3: $\Delta T=qL/k$ and $T_h$ for
$q=5000~\mathrm{W/m^2}$, $L=0.01$~m, $T_c=-40\degC$. Conductivities from Lienhard
Tables A.1 and A.2~\cite{Lienhard}; values from
\pathPone.}
\begin{ruledtabular}
\begin{tabular}{lrlrrp{0.27\textwidth}}
Material & $k$ (W/(m K)) & Source & $\Delta T$ (K) & $T_h$ ($^\circ$C) & Reasonable? \\
\hline
Pure silver (99.99+\%) & 427 & Table A.1, 20\degC & 0.117 & $-39.88$ & Yes: tiny $\Delta T$, far from any limit \\
Pure aluminum & 237 & Table A.1, 20\degC & 0.211 & $-39.79$ & Yes: tiny $\Delta T$, far from any limit \\
Mild steel (0.5\% carbon) & 54 & Table A.1, 20\degC & 0.926 & $-39.07$ & Yes: tiny $\Delta T$, far from any limit \\
Ice & 2.215 & Table A.2, 0\degC & 22.6 & $-17.4$ & Yes: $T_h<0\degC$, ice stays solid \\
Spruce & 0.11 & Table A.2, 20\degC & 455 & 415 & \textbf{No}: wood chars and ignites well below this \\
Insulation (85\% magnesia, \emph{not} MgO) & 0.067 & Table A.2, 38\degC & 746 & 706 & \textbf{No}: far above its service limit \\
Silica aerogel & 0.024 & Table A.2, 0\degC & 2083 & 2043 & \textbf{No}: above the $\mathrm{SiO_2}$ melting point \\
\end{tabular}
\end{ruledtabular}
\end{table*}

\subsection*{Which cases are unreasonable, and why}

\paragraph*{(a) Compare $T_h$ with each material's physical or service limit.}
\begin{itemize}
\item \textbf{Ice:} $T_h\approx-17.4\degC$, still below the 0\degC\ melting point. Ice
remains solid, which is consistent.
\item \textbf{Spruce:} $T_h\approx415\degC$. Wood chars and ignites at roughly
250--300\degC, so spruce cannot exist at this temperature. \textbf{Unreasonable.}
\item \textbf{85\% magnesia:} $T_h\approx706\degC$. The $k=0.067~\WmK$ entry is
\textbf{not} ``MgO''; Lienhard Table A.2 lists ``85\% magnesia'' separately from
``Magnesium oxide, polycrystalline'' ($k=48~\WmK$). 85\% magnesia is a
magnesium-carbonate-based pipe insulation with a documented service limit of
\textbf{up to 315\degC}~\cite{magnesia}. The computed hot side is more than twice that
limit, so this case is \textbf{unreasonable}.
\item \textbf{Aerogel:} $T_h\approx2043\degC$. Silica aerogel is a porous
$\mathrm{SiO_2}$ network. The melting point of $\mathrm{SiO_2}$ is about
\textbf{1710\degC}~\cite{sio2mp}, with reported values of roughly
\textbf{1600--1725\degC}~\cite{sio2range}. 2043\degC\ is above the melting point of the
solid skeleton itself, which is a stronger objection than sintering.
\textbf{Unreasonable.}
\end{itemize}

\paragraph*{(b) Is constant $k$ still credible when $\Delta T$ is hundreds or thousands of kelvin?}
For the three metals and ice, $\Delta T$ is at most a few tens of kelvin, so constant $k$
is fine. For spruce, 85\% magnesia, and aerogel, $\Delta T$ is several hundred to over two
thousand kelvin. Over such ranges $k$ is not constant: it varies with temperature, and for
porous insulators radiation across pores adds a strongly temperature-dependent
contribution. The linear profile derived above no longer holds, so the computed $T_h$ is
only a nominal estimate.

The 85\% magnesia entry makes this concrete: Lienhard Table A.2 gives $k=0.067$ at
38\degC\ and $k=0.08$ at 204\degC. Even within the material's rated range, $k$ rises by
about 20\%. Extrapolating a single constant $k=0.067$ to 706\degC\ is therefore not
defensible.

\paragraph*{(c) Connection to TEC over-temperature protection.}
A thermoelectric cooler uses the Peltier effect to \textbf{pump heat from the cold side to
the hot side}. The element's low $k$ is actually helpful for the cooling function: it
reduces heat conduction back from the hot side to the cold side. The over-temperature
risk is \textbf{not} that the element's own $k$ is low; it lies on the
\textbf{heat-rejection path} on the hot side.

The hot side must reject both the pumped heat and the electrical power through the heat
sink. If the thermal resistance of that path becomes large (a failed fan, poor thermal
contact, a blocked heat sink), the hot-side temperature rise is
\begin{equation*}
\Delta T_{\text{hot side}}\approx\Qdot_{\text{reject}}\,R_{\text{path}}.
\end{equation*}
This has the \textbf{same form} as the slab result $\Delta T=qL/k$: a fixed heat flow
through a large thermal resistance produces a large temperature rise. This is exactly why
Module 4 includes a series thermal switch that opens around 70\degC: it interrupts the
path before the hot side overheats. The slab calculation is a minimal illustration of the
same warning: a large $\Delta T$ can appear when a fixed heat flux is forced through a
high-thermal-resistance path.

%======================================================================
\section*{Part 2. Build the stationary finite-rod model}

\subsection*{2.1 Stationary fin equation}

\paragraph*{Assumptions.}
\begin{enumerate}
\item Steady state: $\partial T/\partial t=0$, so no slice stores energy.
\item No internal heat generation in the rod.
\item $k$, $A$, and $P$ are constant along the rod (a uniform cylinder).
\item $h$ is uniform over the side surface and $T_\infty$ is constant. Radiation is
neglected or linearized into $h$.
\item Each cross section has a single temperature $T(x)$. This requires a small transverse
Biot number, $\Biperp=h(A/P)/k\ll1$. For a circular rod,
$A/P=(\pi d^2/4)/(\pi d)=d/4=r/2$, so $\Biperp=hr/(2k)\ll1$.
\end{enumerate}

\paragraph*{Energy balance.}
Take a slice $[x,\ x+\delta x]$:
\begin{itemize}
\item heat conducted in through the left face (\textbf{Fourier's law}, $q=-k\,dT/dx$,
over the area $A$): $-kA\,\dfrac{dT}{dx}\Big|_x$;
\item heat conducted out through the right face (\textbf{Fourier's law}):
$-kA\,\dfrac{dT}{dx}\Big|_{x+\delta x}$;
\item heat lost from the side (\textbf{Newton's law of cooling}, $q=h(T-T_\infty)$, over
the side area $P\,\delta x$): $hP\,\delta x\,(T-T_\infty)$.
\end{itemize}
\textbf{Conservation of energy} for the slice at steady state (no storage, no
generation): setting in $-$ out $-$ side loss $=0$ gives
\begin{multline}
\left(-kA\frac{dT}{dx}\right)_x-\left(-kA\frac{dT}{dx}\right)_{x+\delta x}\\
-hP\,\delta x\,(T-T_\infty)=0.
\tag{4.28}
\end{multline}
Dividing by $\delta x$ and letting $\delta x\to0$,
\begin{equation*}
kA\,\frac{d^2T}{dx^2}=hP\,(T-T_\infty).
\end{equation*}
Because $T_\infty$ is constant, $d^2T/dx^2=d^2(T-T_\infty)/dx^2$. Therefore
\begin{equation}
\frac{d^2(T-T_\infty)}{dx^2}=m^2\,(T-T_\infty),\qquad m^2=\frac{hP}{kA}.
\tag{4.30}
\end{equation}
For a circular rod, $m^2=4h/(kd)$.

\paragraph*{Unit check.}
Every term in (4.28) has units of W:
\begin{itemize}
\item conduction term: $kA\,dT/dx$ has units
$[\WmK][\mathrm{m^2}][\mathrm{K\,m^{-1}}]=\mathrm{W}$;
\item side-loss term: $hP\,\delta x\,(T-T_\infty)$ has units
$[\WmmK][\mathrm{m}][\mathrm{m}][\mathrm{K}]=\mathrm{W}$.
\end{itemize}
Both sides of $kA\,d^2T/dx^2=hP\,(T-T_\infty)$ are heat flow per unit length,
$\mathrm{W\,m^{-1}}$:
\begin{itemize}
\item left: $[\WmK][\mathrm{m^2}][\mathrm{K\,m^{-2}}]=\mathrm{W\,m^{-1}}$;
\item right: $[\WmmK][\mathrm{m}][\mathrm{K}]=\mathrm{W\,m^{-1}}$.
\end{itemize}
Both sides of (4.30) have units of $\mathrm{K\,m^{-2}}$: the left side is
$\mathrm{K\,m^{-2}}$, and on the right $m^2$ has units
\begin{equation*}
\frac{[\WmmK][\mathrm{m}]}{[\WmK][\mathrm{m^2}]}=\mathrm{m^{-2}},
\end{equation*}
which multiplied by K gives $\mathrm{K\,m^{-2}}$. $1/m$ has units of m and is the decay
length. In an infinite rod, $T-T_\infty$ falls to $1/e$ of its value over each distance
$1/m$.

\subsection*{2.2 Boundary conditions and their physical meaning}

A second-order equation needs two boundary conditions. The course rod uses a Dirichlet
condition at the base and a Neumann condition at the tip.

\paragraph*{Base (Dirichlet):} $T(0)=T_0$. Physically, this fixes the temperature at
$x=0$ to a known value. On the course rod, $x=0$ is the first thermistor hole, and $T_0$
is the temperature that thermistor measures. It is not the TEC setpoint, because the
contact interface and a length of rod lie between the TEC and the first hole, so the two
temperatures differ. Nor is $x=0$ the physical end of the rod. Choosing this interior
plane as the analysis base is valid because the equation holds for any section of rod
from that plane to the tip.

\paragraph*{Tip (Neumann, adiabatic):} $dT/dx|_{x=L}=0$ (4.31b). Physically, no heat
crosses the end face: $\Qdot(L)=-kA\,dT/dx|_L=0$. This is an approximation. The real end
face also loses heat by convection, which is described by the Robin condition (4.31a),
\begin{equation*}
-kA\,\frac{dT}{dx}\Big|_{x=L}=h_LA\,[T(L)-T_\infty].
\end{equation*}
The approximation is reasonable because the ratio of end-face area to side area is
\begin{equation*}
\frac{\pi d^2/4}{\pi dL}=\frac{d}{4L}\ll1,
\end{equation*}
so the end face carries only a small fraction of the total heat loss. Part~3 quantifies
the difference between the two solutions (Example 4.8 and Problem 4.20).

\subsection*{2.3 Finite insulated-tip solution}

\paragraph*{Nondimensionalization.}
Let $\Theta=(T-T_\infty)/(T_0-T_\infty)$ and $\xi=x/L$. Then (4.30) becomes
\begin{equation}
\frac{d^2\Theta}{d\xi^2}=(mL)^2\,\Theta.
\tag{4.34}
\end{equation}
The boundary conditions become $\Theta(0)=1$ and $d\Theta/d\xi|_{\xi=1}=0$. The general
solution and its derivative are
\begin{align*}
\Theta&=C_1e^{mL\xi}+C_2e^{-mL\xi},\\
\frac{d\Theta}{d\xi}&=mL\left(C_1e^{mL\xi}-C_2e^{-mL\xi}\right).
\end{align*}

\paragraph*{Solving for the constants.}
The base condition gives
\begin{equation}
C_1+C_2=1.
\tag{a}
\end{equation}
The tip condition gives
\begin{equation}
mL\left(C_1e^{mL}-C_2e^{-mL}\right)=0\ \Rightarrow\ C_2=C_1e^{2mL}.
\tag{b}
\end{equation}
Substituting (b) into (a),
\begin{align*}
C_1\left(1+e^{2mL}\right)=1\ \Rightarrow\
C_1&=\frac{1}{1+e^{2mL}}=\frac{e^{-mL}}{e^{-mL}+e^{mL}}\\
&=\frac{e^{-mL}}{2\cosh mL}.
\end{align*}
Returning to (b),
\begin{equation*}
C_2=C_1e^{2mL}=\frac{e^{mL}}{2\cosh mL}.
\end{equation*}

\paragraph*{Combining into a hyperbolic cosine.}
\begin{equation}
\begin{split}
\Theta&=\frac{e^{-mL}e^{mL\xi}+e^{mL}e^{-mL\xi}}{2\cosh mL}
=\frac{e^{-mL(1-\xi)}+e^{mL(1-\xi)}}{2\cosh mL}\\
&=\frac{\cosh\big[mL(1-\xi)\big]}{\cosh mL}.
\end{split}
\tag{4.41}
\end{equation}
In dimensional form,
\begin{equation*}
T(x)-T_\infty=(T_0-T_\infty)\,\frac{\cosh\big[m(L-x)\big]}{\cosh mL}.
\end{equation*}

\paragraph*{Verification.}
\begin{itemize}
\item Differential equation:
$\dfrac{d^2\Theta}{d\xi^2}=(mL)^2\dfrac{\cosh[mL(1-\xi)]}{\cosh mL}=(mL)^2\Theta$ $\checkmark$
\item Base condition: $\Theta(0)=\cosh mL/\cosh mL=1$ $\checkmark$
\item Tip condition: $\dfrac{d\Theta}{d\xi}=-mL\,\dfrac{\sinh[mL(1-\xi)]}{\cosh mL}$, which
at $\xi=1$ equals $-mL\,\sinh0/\cosh mL=0$ $\checkmark$
\end{itemize}

\paragraph*{Base heat rate and tip temperature.}
The base heat rate is
\begin{align*}
\Qdot_0&=-kA\frac{dT}{dx}\Big|_0=-\frac{kA(T_0-T_\infty)}{L}\frac{d\Theta}{d\xi}\Big|_0\\
&=kAm\,(T_0-T_\infty)\tanh mL\\
&=\sqrt{hPkA}\,(T_0-T_\infty)\tanh mL.
\end{align*}
Unit check: $\sqrt{hPkA}$ has units
\begin{equation*}
\sqrt{[\WmmK][\mathrm{m}][\WmK][\mathrm{m^2}]}=\mathrm{W\,K^{-1}},
\end{equation*}
and multiplying by K gives W $\checkmark$. The tip temperature is
\begin{equation*}
\Theta(1)=\frac{1}{\cosh mL}.
\end{equation*}

%======================================================================
\section*{Part 3. Provisional finite-length prediction}

\subsection*{3.1 Problem 4.12: the semi-infinite limit}

Problem 4.12 asks for Eq.~(4.50) from the general solution (4.35), using the boundary
conditions $T(x=0)=T_0$ and $T(x=L)=T_\infty$. The far end is held at room temperature;
it is not insulated. With $\theta=T-T_\infty$, these are $\theta(0)=\theta_0$ and
$\theta(L)=0$, and the fin equation $\theta''=m^2\theta$ has the general solution
\begin{equation*}
\theta=C_1e^{mx}+C_2e^{-mx}.
\end{equation*}
The two conditions give $C_1+C_2=\theta_0$ and $C_1e^{mL}+C_2e^{-mL}=0$. The second gives
$C_2=-C_1e^{2mL}$; substituting into the first,
\begin{equation*}
C_1=-\frac{\theta_0\,e^{-mL}}{2\sinh mL},
\qquad
C_2=\frac{\theta_0\,e^{mL}}{2\sinh mL},
\end{equation*}
so
\begin{equation*}
\frac{\theta(x)}{\theta_0}
=\frac{e^{m(L-x)}-e^{-m(L-x)}}{2\sinh mL}
=\frac{\sinh\big[m(L-x)\big]}{\sinh mL}.
\end{equation*}
This satisfies $\theta(0)/\theta_0=1$ and $\theta(L)=0$. For fixed $x$, write it as
\begin{equation}
\frac{\theta(x)}{\theta_0}
=e^{-mx}\,\frac{1-e^{-2m(L-x)}}{1-e^{-2mL}}
\ \xrightarrow{\ L\to\infty\ }\
e^{-mx},
\tag{4.50}
\end{equation}
the semi-infinite profile.

\paragraph*{Physical meaning.}
The far-end condition enters only through the factor $e^{-2m(L-x)}$, which decays
exponentially with the distance from the far end measured in units of $1/m$. A few decay
lengths away from the far end, the temperature no longer depends on what happens there:
the insulated tip (Part~2, $\cosh[m(L-x)]/\cosh mL$) and the room-temperature tip
($\sinh[m(L-x)]/\sinh mL$) both approach $e^{-mx}$. The far-end condition is effectively
forgotten, which is why (4.50) is a limit of the finite solution rather than a new
physical law.

\subsection*{3.2 Example 4.8 and Problem 4.20: the tip-boundary approximation}

For a convecting tip with coefficient $h_L$, the exact solution (Lienhard Eqs.~4.48 and
4.49) is
\begin{equation*}
\Theta=\frac{\cosh\big[mL(1-\xi)\big]+\dfrac{\mathrm{Bi_{tip}}}{mL}\sinh\big[mL(1-\xi)\big]}
{\cosh mL+\dfrac{\mathrm{Bi_{tip}}}{mL}\sinh mL},
\end{equation*}
\begin{equation*}
\frac{\Qdot_0}{\sqrt{hPkA}\,(T_0-T_\infty)}
=\frac{\tanh mL+h_L/(mk)}{1+\big(h_L/(mk)\big)\tanh mL},
\end{equation*}
with $\mathrm{Bi_{tip}}=h_LL/k$ (Lienhard's $\mathrm{Bi_{ax}}$) and
$\mathrm{Bi_{tip}}/(mL)=h_L/(mk)$. Setting $h_L\to0$ recovers the insulated-tip results
(4.41) and (4.44).

\paragraph*{Example 4.8.}
An aluminum rod with $d=2.0$~cm, $k=205~\WmK$, and $L=8.0$~cm protrudes from a 150\degC\
wall into 26\degC\ air with $h=h_L=120~\WmmK$. Then $mL=0.8656$,
$\mathrm{Bi_{ax}}=0.0468$, and $h/(mk)=0.0468/0.8656=0.0541$. The exact tip temperature
is $\Theta_{\mathrm{tip}}=0.6886$ ($T_{\mathrm{tip}}=111.4\degC$), while the insulated tip
gives $\Theta_{\mathrm{tip}}=1/\cosh(0.8656)=0.7150$ ($T_{\mathrm{tip}}=114.7\degC$). The
insulated tip overpredicts the tip excess temperature by about 3.8\%.

\paragraph*{Problem 4.20.}
With $M=\sqrt{hPkA}\,(T_0-T_\infty)$, an insulated tip gives
$\Qdot_{\mathrm{ins}}=M\tanh(mL)$, and a convecting tip gives
\begin{equation*}
\Qdot_{\mathrm{conv}}
=M\,\frac{\sinh(mL)+\frac{h}{mk}\cosh(mL)}{\cosh(mL)+\frac{h}{mk}\sinh(mL)}.
\end{equation*}
Numerically,
\begin{align*}
\frac{\Qdot_{\mathrm{ins}}}{M}&=\tanh(0.8656)=0.6992,\\
\frac{\Qdot_{\mathrm{conv}}}{M}
&=\frac{\sinh(0.8656)+0.0541\cosh(0.8656)}{\cosh(0.8656)+0.0541\sinh(0.8656)}=0.7258,
\end{align*}
so the insulated-tip approximation underpredicts the exact heat flow by
\begin{equation*}
\frac{\Qdot_{\mathrm{conv}}-\Qdot_{\mathrm{ins}}}{\Qdot_{\mathrm{conv}}}
=\frac{0.7258-0.6992}{0.7258}=0.0367,
\end{equation*}
or \textbf{3.67\%}, with the approximate heat flow slightly lower than the exact value.

\paragraph*{Course rod.}
The same comparison with $h_L=h=10~\WmmK$, $k=167~\WmK$, and the inputs of 3.3 gives
$mL=1.168$, $\mathrm{Bi_{tip}}=hL/k=0.0161$, and $h/(mk)=0.0138$, four times smaller than
in Example 4.8 (Table~\ref{tab:tip}). The insulated-tip boundary condition therefore
changes the predicted course-rod profile by about 1\% at most, at the tip itself.

\begin{table}[b]
\caption{\label{tab:tip}Insulated versus convecting tip for the course rod
($h_L=h=10~\WmmK$, $k=167~\WmK$, $L=26.90$~cm).}
\begin{ruledtabular}
\begin{tabular}{lrrr}
Quantity & Insulated & Convecting & Difference \\
\hline
$\Theta_{\mathrm{tip}}$ & 0.5671 & 0.5607 & insulated 1.1\% high \\
$\Qdot_0/M$ & 0.8236 & 0.8280 & insulated 0.53\% low \\
\end{tabular}
\end{ruledtabular}
\end{table}

\subsection*{3.3 Inputs}

The inputs are listed in Table~\ref{tab:inputs}.

\begin{table*}[t]
\caption{\label{tab:inputs}Inputs for the provisional predictions of Parts 3 and 4.}
\begin{ruledtabular}
\begin{tabular}{lp{0.33\textwidth}p{0.4\textwidth}}
Input & Value & Source \\
\hline
Rod diameter $d$ & $12.70\pm0.02$ mm & measured with digital calipers (0.01 mm resolution) \\
Overall rod length & $28.66\pm0.05$ cm & measured with a ruler (1 mm divisions) \\
Hole diameter & 2.50--2.55 mm (mean 2.525 mm) & measured with calipers \\
First-hole offset $s$ & $16.38~\mathrm{mm}+\tfrac12(2.525~\mathrm{mm})=17.64$ mm & calipers, from the TEC end of the rod to the near edge of the first hole, plus half the hole diameter to reach its center \\
Analyzed length $L$ & $28.66~\mathrm{cm}-1.764~\mathrm{cm}=26.90\pm0.05$ cm & first-hole center to the tip \\
Hole positions $x$ & 0, 1.2, 2.4, 3.7, 5.0, 7.6, 10.2, 15.3, 20.4 cm & measured from the first hole \\
Planned sensors & T1--T5 at $x=0$, 5.0, 10.2, 15.3, 20.4 cm & chosen below \\
Position uncertainty $\delta x$ & $\pm1.3$ mm & half the hole diameter: a sensor can sit anywhere in its hole \\
Side coefficient $h$ & $10~\WmmK$ & Part III planning value for still room air~\cite{handout}; not a measurement \\
Conductivity $k$ & 167 W/(m K); range 166--172 W/(m K) & 6061-T6, Lienhard Table A.1: 20\degC\ value; range from the 0\degC\ and 100\degC\ values \\
Error criterion & $\epsilon_T<5\%$ at every sensor used & Part III handout's suggested starting criterion \\
\end{tabular}
\end{ruledtabular}
\end{table*}

\paragraph*{Planned sensor positions.}
T1--T5 go in the holes at $x=0$, 5.0, 10.2, 15.3, and 20.4 cm. The reasons, using the
predictions of 3.4 for $L=26.90$ cm and $h=10~\WmmK$:
\begin{enumerate}
\item T1 must be at $x=0$: the first hole is the analysis base defined in the handout.
\item The farthest hole, 20.4 cm, is essential. The decay length $1/m\approx23$ cm is
longer than the measured region, and even there the predicted $\Theta$ is only about 0.59,
so the whole available range is needed to measure a large enough temperature decay.
\item The points are spread evenly on a plot of $\ln\theta$ against $x$, as the handout
asks, rather than clustered at one end: the spacings are all about 5 cm, and the
predicted $\ln\Theta$ values are 0, $-0.17$, $-0.33$, $-0.45$, and $-0.53$.
\item The set tests the difference between the finite and semi-infinite models: 5.0 cm is
the last hole that meets the 5\% criterion ($\epsilon_T=4.6\%$), while at 10.2, 15.3, and
20.4 cm $\epsilon_T$ is 11\%, 20\%, and 30\%, so the two fits should differ visibly in
their residuals.
\item The holes at 1.2 and 2.4 cm are not used: their 1.2 cm spacing is only about ten
times the $\pm1.3$ mm position uncertainty, and neighboring points would differ little in
temperature, adding little information.
\item The farthest sensor still has a large signal: with the base about 25 K above room
temperature, the predicted excess temperature at 20.4 cm is about 15 K, far above the
noise.
\end{enumerate}

\subsection*{3.4 $m$, $mL$, and $\epsilon_T(x)$ at each planned sensor}

For the circular rod, $P/A=4/d$, so
\begin{align*}
m=\sqrt{\frac{4h}{kd}}
&=\sqrt{\frac{4(10~\WmmK)}{(167~\WmK)(0.01270~\mathrm{m})}}\\
&=4.34~\mathrm{m^{-1}}.
\end{align*}
This gives $m=4.343~\mathrm{m^{-1}}$, $1/m=23.0$ cm, $mL=1.168$, and
$\mathrm{Bi}_\parallel=(mL)^2=1.364$. Over $k=166$--172 W/(m K), $m=4.279$--4.356
$\mathrm{m^{-1}}$ and $mL=1.151$--1.172; this range changes the results below by much
less than their rounding.

\paragraph*{Error formula.}
Dividing the insulated-tip solution by the semi-infinite one,
\begin{equation*}
\frac{\Theta_{\mathrm{finite}}}{\Theta_\infty}
=\frac{\cosh[m(L-x)]}{\cosh(mL)\,e^{-mx}}
=\frac{1+e^{-2m(L-x)}}{1+e^{-2mL}}\ \ge\ 1,
\end{equation*}
so the insulated tip keeps the rod warmer than the semi-infinite model predicts, and the
difference is largest near the tip. The relative error of the semi-infinite model is
therefore
\begin{equation*}
\epsilon_T(x)
=\frac{\Theta_{\mathrm{finite}}-\Theta_\infty}{\Theta_{\mathrm{finite}}}
=1-\frac{1+e^{-2mL}}{1+e^{-2m(L-x)}},
\end{equation*}
which increases monotonically from 0 at the base. Setting $\epsilon_T=5\%$ gives the
farthest position at which the criterion holds, $x^{\ast}=5.4$ cm. Table~\ref{tab:holes}
gives $\epsilon_T$ at every hole and Table~\ref{tab:sensors} at the planned sensors.
Because $h$ is only a planning value, Table~\ref{tab:hsens} varies it at $k=167$ W/(m K).
Values are from \path{python/analysis/a4_part3_finite_length.py}.

\begin{table*}[t]
\caption{\label{tab:holes}Finite insulated-tip and semi-infinite profiles at every hole
for $h=10~\WmmK$, $k=167~\WmK$, $L=26.90$ cm.}
\begin{ruledtabular}
\begin{tabular}{rrrrrrc}
Hole $x$ (cm) & $L-x$ (cm) & $\Theta_{\mathrm{finite}}$ & $\Theta_\infty$ & $\epsilon_T$ & $\epsilon_T$ for $x\pm\delta x$ & $\epsilon_T<5\%$? \\
\hline
0.0 & 26.90 & 1.0000 & 1.0000 & 0.00\% & --- & yes \\
1.2 & 25.70 & 0.9584 & 0.9492 & 0.96\% & 0.85--1.07\% & yes \\
2.4 & 24.50 & 0.9194 & 0.9010 & 2.00\% & 1.89--2.12\% & yes \\
3.7 & 23.20 & 0.8800 & 0.8516 & 3.23\% & 3.11--3.36\% & yes \\
5.0 & 21.90 & 0.8434 & 0.8048 & 4.58\% & 4.44--4.71\% & yes \\
7.6 & 19.30 & 0.7782 & 0.7189 & 7.62\% & 7.46--7.78\% & no \\
10.2 & 16.70 & 0.7228 & 0.6421 & 11.16\% & 10.98--11.35\% & no \\
15.3 & 11.60 & 0.6406 & 0.5146 & 19.67\% & 19.44--19.91\% & no \\
20.4 & 6.50 & 0.5898 & 0.4123 & 30.09\% & 29.82--30.37\% & no \\
\end{tabular}
\end{ruledtabular}
\end{table*}

\begin{table}[t]
\caption{\label{tab:sensors}Planned sensors, with an illustrative base excess temperature
$\theta_0=25$ K.}
\begin{ruledtabular}
\begin{tabular}{lrrrrr}
Sensor & $x$ (cm) & $\Theta_{\mathrm{finite}}$ & $\ln\Theta_{\mathrm{finite}}$ & $\theta$ (K) & $\epsilon_T$ \\
\hline
T1 & 0.0 & 1.0000 & 0.000 & 25.0 & 0.00\% \\
T2 & 5.0 & 0.8434 & $-0.170$ & 21.1 & 4.58\% \\
T3 & 10.2 & 0.7228 & $-0.325$ & 18.1 & 11.16\% \\
T4 & 15.3 & 0.6406 & $-0.445$ & 16.0 & 19.67\% \\
T5 & 20.4 & 0.5898 & $-0.528$ & 14.7 & 30.09\% \\
\end{tabular}
\end{ruledtabular}
\end{table}

\begin{table}[t]
\caption{\label{tab:hsens}Sensitivity to the planning value of $h$ at $k=167~\WmK$.}
\begin{ruledtabular}
\begin{tabular}{rrrr}
$h$ ($\WmmK$) & $1/m$ (cm) & $mL$ & $x^{\ast}$ (cm) \\
\hline
5 & 32.6 & 0.826 & 4.6 \\
10 & 23.0 & 1.168 & 5.4 \\
15 & 18.8 & 1.431 & 6.4 \\
20 & 16.3 & 1.652 & 7.4 \\
\end{tabular}
\end{ruledtabular}
\end{table}

\subsection*{3.5 Prediction: length}

With the criterion $\epsilon_T<5\%$ at every sensor used, the semi-infinite model is
adequate at T1 ($x=0$) and T2 (5.0 cm, $\epsilon_T=4.6\%$, still below 5\% at the edge of
the $\pm1.3$ mm position uncertainty), but not at T3, T4, and T5 (10.2, 15.3, and
20.4 cm), where $\epsilon_T$ is 11\%, 20\%, and 30\%. Among all the holes, only those at
$x\le5.0$ cm meet the criterion ($x^{\ast}=5.4$ cm). The whole rod is far from the
``infinitely long'' regime: $mL=1.17$, compared with Lienhard's guideline $mL\gtrsim5$.

Three of the five planned sensors lie beyond $x^{\ast}$, so I expect the course rod
\textbf{not} to be adequately semi-infinite over the measured region. The stationary
profile should be analyzed with the finite insulated-tip model; the semi-infinite fit is a
comparison that should show a systematic residual at the far sensors. The tip boundary
condition itself matters much less: replacing the insulated tip by a convecting one
changes the prediction by about 1\% (3.2).

This prediction depends on the planned sensor positions and $L=26.90$ cm in 3.3,
$h=10~\WmmK$, $k=166$--172 W/(m K), and the 5\% criterion. It is not sensitive to $k$.
For any $h$ from 5 to 20 $\WmmK$, $x^{\ast}$ stays between 4.6 and 7.4 cm, so T3, T4, and
T5 fail in every case. Only T2 at 5.0 cm is marginal: it passes for $h\ge10~\WmmK$ but
fails if $h$ is as low as 5 $\WmmK$.

%======================================================================
\section*{Part 4. Provisional one-dimensional prediction}

\subsection*{4.1 Transverse Biot number}

\paragraph*{Definition and meaning.}
The transverse (radial) Biot number is defined as (Lienhard Eq.~4.27)
\begin{equation*}
\Biperp=\frac{h(A/P)}{k},
\end{equation*}
where $A/P$, the cross-sectional area divided by the perimeter, is the characteristic
length for conduction across the section. It represents the ratio of the internal radial
conduction resistance to the external surface convection resistance:
\begin{equation*}
\Biperp\sim\frac{R_{\mathrm{radial\ conduction}}}{R_{\mathrm{surface\ convection}}}.
\end{equation*}

\paragraph*{Form for a cylindrical rod.}
For a rod of diameter $d$,
\begin{equation*}
\frac{A}{P}=\frac{\pi d^2/4}{\pi d}=\frac{d}{4},
\qquad\text{so}\qquad
\boxed{\Biperp=\frac{hd}{4k}}.
\end{equation*}

\paragraph*{Numerical value and unit check.}
Using $h=10~\WmmK$, $d=0.01270$ m, and $k=167~\WmK$,
\begin{equation*}
\Biperp=\frac{(10)(0.01270)}{4(167)}=1.90\times10^{-4}.
\end{equation*}
If $k$ is taken from Lienhard's range of 166--172 W/(m K), then
$\Biperp=1.85\times10^{-4}$ to $1.91\times10^{-4}$. Units:
\begin{equation*}
\frac{[\WmmK][\mathrm{m}]}{[\WmK]}=1,
\end{equation*}
so $\Biperp$ is dimensionless.

\paragraph*{Conversion to a temperature difference.}
The temperature difference between the center of the cross section and its surface is
approximately
\begin{equation*}
\Delta T_{\mathrm{center-surface}}\approx\Biperp\,\theta,
\end{equation*}
where $\theta$ is the temperature difference between the rod surface and the ambient air.
With the illustrative $\theta\approx25$ K of 3.4,
\begin{equation*}
\Delta T\approx(2\times10^{-4})(25~\mathrm{K})\approx0.005~\mathrm{K},
\end{equation*}
about ten times smaller than the scatter of single thermistor readings (standard deviation
0.03--0.05\degC\ in the stationary interval of the rod runs).

\paragraph*{Inverse calculation: what $h$ would break the radial lumped assumption?}
Setting $\Biperp=0.1$,
\begin{equation*}
h=\frac{0.4k}{d}=\frac{0.4(167)}{0.0127}\approx5300~\WmmK.
\end{equation*}
Lienhard Table 1.1 (p.~21) lists 4.2--29 $\WmmK$ for natural convection in air and
45--80 $\WmmK$ for forced convection of air. The value 5300 $\WmmK$ is about 500 times the
planning value and still about 65 times the largest forced-air value. Under the
experimental conditions, the assumption of radial temperature uniformity is therefore
very robust.

\subsection*{4.2 Prediction: radius}

\paragraph*{Criterion.}
Lienhard (p.~24): the lumped-capacity solution is normally accurate within about 3\% if
$\mathrm{Bi}<0.1$.

\paragraph*{Conclusion.}
For this rod, $\Biperp\approx2\times10^{-4}\ll0.1$. Therefore each cross section can be
represented by a single temperature $T(x)$. This conclusion depends on $h=10~\WmmK$,
$k=166$--172 W/(m K), and $d=12.70$ mm.

\paragraph*{Distinction from 3.4 and 3.5.}
$\Biperp$ answers the approximation question in the \textbf{radial} direction;
$\mathrm{Bi}_\parallel=(mL)^2$ in 3.4 and the prediction in 3.5 answer it in the
\textbf{length} direction. They are independent: for this rod $\Biperp$ is very small,
but $(mL)^2=1.36$ is not. The temperature is uniform across each section, yet it still
decays noticeably along the rod.

\paragraph*{Status of this prediction.}
This is a prediction made \emph{before} heating. After $m$ is fitted, the
self-consistency check
\begin{equation*}
\Biperp=\left(\frac{md}{4}\right)^2
\end{equation*}
will confirm whether the experimentally extracted $m$ is consistent with a small
$\Biperp$.

\subsection*{4.3 Application: heat sink in still and forced air}

\paragraph*{Larger $h$ in forced air.}
Forced airflow continuously sweeps away the heated air near the surface, thinning the
thermal boundary layer, and thus significantly increases the convective heat transfer
coefficient $h$. Lienhard Table 1.1 gives 4.2--29 $\WmmK$ for natural convection in air
and 45--80 $\WmmK$ for forced convection of air (illustrative values; the actual $h$
depends on velocity and geometry).

\paragraph*{Relation between the decay length $1/m$ and $h$.}
From $m=\sqrt{hP/(kA)}$,
\begin{equation*}
\frac{1}{m}=\sqrt{\frac{kA}{hP}}=\sqrt{\frac{kd}{4h}}\propto h^{-1/2}.
\end{equation*}
With $k=167~\WmK$ and $d=0.01270$ m, $1/m=32.6$, 23.0, 18.8, and 16.3 cm for $h=5$, 10,
15, and 20 $\WmmK$ (Table~\ref{tab:hsens}). When $h$ is quadrupled, $1/m$ is halved. The
forced-air heat sink therefore has the shorter decay length.

\paragraph*{Physical meaning.}
A shorter decay length means that heat leaves through the sides before it can travel far
along the fin, so the far end of the fin is closer to the ambient air temperature. The
total heat dissipation,
\begin{equation*}
\Qdot_0=\sqrt{hPkA}\,\theta_0\tanh(mL),
\end{equation*}
still increases with $h$. $\Biperp\propto h$ also increases with $h$, but it remains
$\ll1$ in air.

%======================================================================
\section*{Git checkpoint and AI use}

This PDF is built from commit \texttt{\commit} of \repourl; the same text is in
\path{docs/assessments/a4_finite_length_biot.md}, and the LaTeX source in
\path{docs/assessments/a4_latex/A4_Yin.tex}.

\begin{itemize}
\item \textbf{Part 1:} The derivation, the reasonableness discussion, and the choice of
references are my own. Claude Code (an AI coding assistant) wrote
\path{python/analysis/a4_problem_1_3.py}, which computes Table~\ref{tab:p13} from the
Lienhard Table A.1 and A.2 values, reviewed my drafts, and formatted the references.
\item \textbf{Part 2:} The calculations are my own. Claude Code organized and formatted
this part: it wrote a first draft for comparison, reviewed my version, and fixed its
LaTeX formatting.
\item \textbf{Part 3:} The Problem 4.12 derivation, the Example 4.8 and Problem 4.20
calculations, and the choice of sensor holes with its reasons are my own. The course-rod
numbers in 3.2--3.4 ($m$, $mL$, $\epsilon_T$ at each hole, $x^{\ast}$, the
planned-sensor table, the tip comparison, and the $h$ sensitivity) were computed by
\path{python/analysis/a4_part3_finite_length.py}, which Claude Code wrote from my
measurements; I checked $m$, $mL$, and $\epsilon_T$ at 7.6 cm by hand. Claude Code also
organized and formatted Part 3: it made the tables, translated my notes into English, and
drafted the text of 3.4 and 3.5.
\item \textbf{Part 4:} The text and calculations are my own. Claude Code checked the
numbers and the Lienhard citations against the textbook, corrected the page of the 3\%
statement (p.~24), the ``two orders of magnitude'' comparison, and the reading-noise value
(now the measured scatter), and converted the math to GitHub format.
\item \textbf{This PDF:} Claude Code converted the repository Markdown to this REVTeX
document and drafted the abstract.
\end{itemize}

\begin{thebibliography}{6}
\bibitem{Lienhard}
J.~H. Lienhard~V and J.~H. Lienhard~IV, \emph{A Heat Transfer Textbook}, 6th ed.
(Phlogiston Press, Cambridge, MA, 2024).

\bibitem{handout}
Phys 39, Brandeis University, Fall 2026, \emph{Module 7, Part III: Rod Instrument and
Stationary Fin}, \url{https://sethfraden.github.io/Phys39F26-course/labs/lab-08/}.

\bibitem{magnesia}
Asbestos Safety and Eradication Agency (Australia), ``85\% Magnesia insulation,''
\url{https://products.asbestossafety.gov.au/insulation/85-magnesia-insulation/}.

\bibitem{sio2mp}
WebElements, ``Silicon dioxide,''
\url{https://www.webelements.com/compounds/silicon/silicon_dioxide.html}.

\bibitem{sio2range}
Wikipedia, ``Silicon dioxide'' (archived 13 November 2013),
\url{https://eresources.nlb.gov.sg/webarchives/wayback/20131113210036/https://en.wikipedia.org/wiki/Silica}.
\end{thebibliography}

\end{document}
